%%%%%%%%%%%%%%%%%%%%%%%%%%%%%%%%%%%%%%%%%%%%%%%%%%%%%%%%%%%%%%%%%%%%%%%%%%%%%%%%
%2345678901234567890123456789012345678901234567890123456789012345678901234567890
%        1         2         3         4         5         6         7         8

\documentclass[a4, 10 pt, conference]{ieeeconf}

\IEEEoverridecommandlockouts
\overrideIEEEmargins

\usepackage{multicol}
\usepackage{tcolorbox}
\usepackage{cuted,tcolorbox,lipsum}
\usepackage{xcolor}
\usepackage{graphicx}
\usepackage{booktabs}
\usepackage{multirow}
\usepackage{amsmath}

\title{\LARGE \bf
Introduction to Machine Learning (SS 2026)\\ Programming Project
\vspace{-3em}
}

\begin{document}

\maketitle
\vspace{-3em}
\thispagestyle{empty}
\pagestyle{empty}

\begin{strip}
\begin{tcolorbox}[
size=tight,
colback=white,
boxrule=0.2mm,
left=3mm,right=3mm, top=3mm, bottom=1mm
]
{\begin{multicols}{2}

\textbf{Author 1}           \\
Last name: Fischbach        \\
First name: Matti Sebastian \\
Matrikel Nr.: 12404020

\columnbreak

\textbf{Author 2} \\
Last name: Musch  \\
First name: Eric  \\
Matrikel Nr.: 12320608

\end{multicols}}
\end{tcolorbox}
\end{strip}

%%%%%%%%%%%%%%%%%%%%%%%%%%%%%%%%%%%%%%%%%%%%%%%%%%%%%%%%%%%%%%%%%%%%%%%%%%%%%%%%

\section{Introduction}
\label{sec:intro}

We address a binary classification problem: detecting fraudulent financial
transactions. The dataset contains 227{,}845 transactions described by 31
features: \texttt{Time}, \texttt{Amount}, and 28 anonymized principal
components (\texttt{Feature0}–\texttt{Feature27}) produced by a prior PCA
transformation. The target variable \texttt{Class} encodes legitimate (0) or
fraudulent (1) transactions.

The dataset is severely imbalanced: only 394 of 227{,}845 transactions are
fraud (0.17\%). This makes accuracy a misleading metric, since a trivial
all-negative classifier achieves 99.83\% accuracy while catching zero fraud.
We therefore report Area Under the Precision–Recall Curve (AUPRC) as our
primary metric, which directly captures the precision–recall trade-off on the
minority class and is not inflated by the large number of true negatives.
ROC-AUC is used as a secondary metric for leaderboard/test-dataset evaluation.

\section{Implementation / ML Process}
\label{sec:methods}

\subsection{Preprocessing}

We applied an 80/20 stratified train-validation split to preserve the 0.17\%
fraud rate in both sets. The \texttt{Time} feature was dropped after empirical
analysis showed it provides negligible predictive signal (AUPRC delta $<$0.001).
\texttt{Amount} was scaled with \texttt{RobustScaler} fitted on the training set
only to prevent data leakage; this is required for Logistic Regression but
omitted for Random Forest as tree models are invariant to monotone feature
transformations. Class imbalance was handled via \texttt{class\_weight=None}
on the LR model and \texttt{class\_weight=balanced} on the RF model.

\subsection{Method 1: Logistic Regression}

Logistic Regression serves as our interpretable linear baseline. We ran
\texttt{GridSearchCV} (3-fold CV, scored on \texttt{average\_precision}) over $C
\in \{0.1, 1, 10\}$ and both L1 and L2 regularization, with
\texttt{class\_weight} $\in \{\texttt{balanced},\texttt{None}\}$. The best
configuration was L1 regularization (\texttt{saga} solver, implemented via
\texttt{elasticnet} with \texttt{l1\_ratio=1}), $C=0.1$,
\texttt{class\_weight=None}.

\subsection{Method 2: Random Forest}

The Random Forest classifier uses the Decision Tree as its underlying
architecture. In a decision tree we have Internal Nodes (split feature $x_d$
with threshold $\theta$) and Leaves which in our case are class labels (our
classification). The splitting is done by reducing the \textbf{impurity} at each
node. In the case of the \texttt{scikit-learn} library used, this is done via
Gini impurity. The random forest uses the decision tree architecture and uses
\textbf{bagging} (or bootstrap aggregation) to decorrelate the individual
decision trees and reduce variance. We do this by sampling $K$ datasets with
replacement from the training dataset, then train each decision tree
independently on a different bootstrapped dataset and average the predicted
class probabilities across all of our $K$ trees.

Averaging over 500 trees reduces variance substantially without overfitting,
which is a key advantage over single decision trees on this high-dimensional
dataset. Unlike Logistic Regression, Random Forest can capture non-linear
interactions between the anonymized PCA features.

Hyperparameters were selected through five rounds of \texttt{RandomizedSearchCV}
and \texttt{GridSearchCV} (3-fold CV, \texttt{average\_precision} scoring),
progressively narrowing the search space. The key findings were:
\texttt{max\_depth=None} (unlimited depth) consistently outperformed depth-limited
trees because the 500-tree averaging prevents overfitting;
\texttt{max\_features=log2} produced more diverse trees than \texttt{sqrt} at
high tree counts;
and \texttt{max\_samples=None} (full bootstrap sample) was confirmed optimal.
Final parameters: \texttt{n\_estimators=500}, \texttt{max\_depth=None},
\texttt{max\_features=log2}, \texttt{class\_weight=balanced}.

\subsection{Decision Threshold}

For both models, the default classification threshold of 0.5 is suboptimal under
heavy class imbalance. We swept the Precision–Recall curve and selected the
threshold that maximizes F1-score on the validation set. This threshold is saved
alongside the model artifact.

\section{Results}
\label{sec:results}

Table~\ref{tab:metrics} summarises validation set performance for all four model
variants (simple = default hyperparameters; tuned = post-search best params).

\begin{table}[h]
\centering
\caption{Validation Set Performance}
\label{tab:metrics}
\begin{tabular}{lrrrr}
\toprule
Model & AUPRC & ROC-AUC & F1 & Recall \\
\midrule
LR Simple    & 0.787 & 0.981 & 0.857 & 0.835 \\
LR Tuned     & 0.807 & 0.983 & 0.827 & 0.848 \\
RF Simple    & 0.887 & 0.955 & 0.899 & 0.848 \\
RF Tuned     & 0.899 & 0.966 & 0.903 & 0.886 \\
\bottomrule
\end{tabular}
\end{table}

Figure~\ref{fig:confusion} shows the confusion matrix of the tuned Logistic
Regression and Random Forest models at their F1-optimal threshold
(\texttt{LR\_threshold$\approx$0.07}; \texttt{RF\_threshold$\approx$0.29}). Our
best model (Random Forest) catches 70 of the 79 fraud cases (recall 0.886) while
raising only 6 false alarms out of 45{,}490 legitimate transactions (precision
0.921).

\begin{figure}[h]
  \centering
  \includegraphics[width=\columnwidth]{figures/compare_confusion_matrices}
  \caption{Confusion matrix for LR and RF tuned.}
  \label{fig:confusion}
\end{figure}

Figure~\ref{fig:pr} shows Precision–Recall curves for all four models.
The tuned RF achieves the highest AUPRC and dominates across all recall levels.
Both simple models are already competitive, with tuning providing a consistent
improvement of approximately +0.02 (LR) and +0.01 (RF).

\begin{figure}[h]
  \centering
  \includegraphics[width=\columnwidth]{figures/compare_pr_curves}
  \caption{Precision–Recall curves for LR and RF (simple and tuned).
           Dots mark the F1-optimal operating threshold per model.}
  \label{fig:pr}
\end{figure}

Figure~\ref{fig:importance} shows the top-20 feature importances for the tuned
RF model. \texttt{Feature13} is clearly the most important feature, followed
by three features (\texttt{Feature3}, \texttt{Feature9}, \texttt{Feature4})
with similar, substantially lower importance.

\begin{figure}[h]
  \centering
  \includegraphics[width=\columnwidth]{figures/rf_feature_importance}
  \caption{RF feature importances (mean impurity decrease, tuned model).}
  \label{fig:importance}
\end{figure}

\section{Discussion}
\label{sec:discuss}

Random Forest outperforms Logistic Regression substantially in AUPRC (+0.09),
which we attribute to non-linear interactions among the PCA features that a
linear model cannot capture. The high ROC-AUC (0.966 val / 0.972 test) confirms
the model generalises well to unseen data.

\textbf{What was tried and discarded:}
(1) Including \texttt{Time} as a feature (AUPRC delta was $<$0.001), so it was
dropped to simplify the pipeline.
(2) A Neural Network was initially considered as the second method but replaced
by Random Forest, which requires no architecture design and achieves better
out-of-the-box performance on tabular data.
(3) During tuning, accidentally omitting \texttt{class\_weight=balanced}
produced the fewest false positives (3) but lower recall than the balanced
variant (68/79 fraud caught). This further confirms that balanced weighting is
critical for this highly imbalanced dataset.
(4) Tuned LR achieves higher AUPRC (0.807 vs.\ 0.787) but lower F1 (0.827 vs.\
0.857) than simple LR. Stronger regularization ($C=0.1$, L1) improves
probability calibration across all thresholds (higher AUPRC) but constrains the
peak F1 attainable at the optimal operating point.

\textbf{Potential improvements:}
Gradient boosting methods (XGBoost, LightGBM) may further improve AUPRC.
Probability calibration could refine threshold selection.

\section{Conclusion}
\label{sec:con}

We implemented two fraud detection classifiers, Logistic Regression and
Random Forest, on a severely imbalanced financial transactions dataset
(0.17\% fraud). The tuned Random Forest achieves an AUPRC of 0.899 on the
local validation set and a ROC-AUC of \textbf{0.9716} on the hidden test set,
demonstrating strong generalisation. Our main takeaways are:
(1) AUPRC is the appropriate metric for heavily imbalanced binary classification;
(2) For RF, \texttt{class\_weight=balanced} is essential for recall on the
minority class; LR's grid search found \texttt{class\_weight=None} optimal for
AUPRC.
(3) Random Forest with unlimited tree depth and 500 estimators significantly
outperforms a tuned linear baseline, capturing non-linear structure in the
anonymised PCA features without requiring feature scaling or resampling.

\end{document}
